% Figure: a vertical pole held by a single guy wire, loaded at the top by
% a horizontal pull. The guy tension and the ground anchor carry the load;
% the base sees a large vertical thrust.
\begin{figure}[htbp]
\centering
\begin{tikzpicture}[
  pole/.style={draw=engblack, line width=2.0pt},
  wire/.style={draw=engblue, line width=1.2pt},
  force/.style={draw=engred, -{Stealth}, very thick},
  reac/.style={draw=englime!70!engblack, -{Stealth}, very thick},
  ground/.style={draw=enggray, thin},
  scale=1.0,
]
% ground line
\draw[ground] (-0.6,0) -- (5.2,0);
\foreach \x in {-0.4,0.0,0.4,0.8,1.2,1.6,2.0,2.4,2.8,3.2,3.6,4.0,4.4,4.8} {
  \draw[ground] (\x,0) -- ++(-0.18,-0.18);
}
% pole base at A, top at B
\coordinate (A) at (0,0);
\coordinate (B) at (0,4.0);
% anchor on the ground at C, to the right
\coordinate (C) at (3.0,0);
\draw[pole] (A) -- (B);
\draw[wire] (B) -- (C) node[midway, anchor=south west, font=\scriptsize, engblue] {$T$};
% base and anchor markers
\fill[engblack] (A) circle (0.07);
\node[font=\scriptsize, anchor=north east] at (A) {$A$};
\fill[engblack] (C) circle (0.07);
\node[font=\scriptsize, anchor=north] at (C) {$C$};
\node[font=\scriptsize, anchor=east] at (B) {$B$};
% applied horizontal load P at the top, pulling left
\draw[force] (B) -- ++(-1.4,0) node[anchor=east, font=\small, engred] {$P$};
% angle of the guy from the pole (vertical)
\draw[draw=enggray, thin] (0,3.2) arc (-90:-53:0.8);
\node[font=\scriptsize] at (0.55,3.05) {$\psi$};
% base reaction: large vertical thrust plus horizontal
\draw[reac] (A) -- ++(0,-1.2)
  node[anchor=north, font=\scriptsize, englime!60!engblack] {$A_y$ (down on pole)};
\end{tikzpicture}
\caption[A guyed pole loaded at the top]{A vertical pole pinned at its
base $A$ is held against a horizontal pull $P$ at the top $B$ by a single
guy wire running to a ground anchor $C$. The guy makes an angle $\psi$
with the pole. The guy can pull only along its own line, so its horizontal
component must balance $P$ while its vertical component presses the pole
down onto its base. A shallow guy (large $\psi$, anchor set far out)
balances $P$ cheaply in tension but with a modest downthrust; a steep guy
(small $\psi$, anchor close to the base) needs a large tension to find any
horizontal component, and drives a correspondingly large vertical thrust
into the base. The base, not the wire, often governs.}
\label{fig:vol03ch01:guy-wire}
\end{figure}
